% fig-nir.tex -- NIR as a discrimination score. Goes with the metric definition
% in Sections/04a-ruler.tex, and is \input from there.
%
% No data dependency: this is a schematic. Its one numeral, NIR = 6/8, is read
% off the drawing itself -- six of the eight hollow circles lie left of the
% filled dot -- and agrees with the definition and the tie-term paragraph in
% 04a-ruler.tex. The sibling fig-nir.json records that check.
%
% This is the v2 COPY of thesis/figs/fig-nir.tex. The v1 file is frozen and was
% not touched. Two changes, and nothing else.
%
% (1) MEASURE. The v1 drawing was 302.79pt wide against v2's 404.03pt (142mm)
%     block -- 75% of the measure, so it sat as a small island with 36mm of
%     white at its right. It is redrawn at the full measure: every x-coordinate
%     is v1's times 1.598, and the marks that ARE the drawing (dot radii, band
%     height, brace amplitude, arrowhead) scale with them, so the figure is
%     enlarged rather than stretched. Type is NOT scaled: \scriptsize and \small
%     stay at the document's own sizes. That is the whole point -- v1's figures
%     were scaled by LaTeX and their type shrank with them. The vertical rhythm
%     of the three reading rules opens from 0.34cm to 0.42cm, proper leading for
%     8pt type rather than the 9.7pt v1 had.
%
%     The explicit bounding box below belongs to this change. A TikZ bounding
%     box includes each node's invisible inner sep, so a picture whose BOX is
%     flush with the measure has its INK sitting 3.5mm inside it at both edges,
%     and the figure reads as not quite reaching the margin. Declaring the box
%     to be the ink fixes that. Measured on the built page: ink 141.82mm inside
%     the 142.00mm block, 0.31pt clear of the left margin and 0.21pt of the
%     right. Nothing is scaled in LaTeX -- no \resizebox, no scale= key, no
%     width= key -- and the document builds with zero overfull boxes.
%
% (2) LEADER. "own truth" floated above the filled dot with a hollow circle
%     1.4cm to its right, and at this width a reader can take it to be labelling
%     that hollow circle. A 0.35cm leader in quietink now ties the label to the
%     dot it names. A plain line, not an arrow: it identifies, it does not
%     assert a direction.
%
% NOT changed: the nine circles and their irregular spacing, NIR = 6/8, all
% three reading rules including the half-tie rule (\S4.7 turns on the half-tie
% property, so its in-figure statement has to survive), and the caption.
\begin{figure}[htbp]
% MARGINALIA (06-figures.md): wider than the 124mm text column, so it
% sits in fullwidth; the caption below spans the text column only.
\begin{fullwidth}\raggedright
\begin{adjustbox}{max width=\linewidth}
\begin{tikzpicture}[
  x=1cm, y=1cm, font=\small,
  lbl/.style={font=\scriptsize, text=ink!70},
  strong/.style={font=\scriptsize, text=ink!90},
  own/.style={circle, fill=ink!85, inner sep=3.2pt},
  other/.style={circle, draw=ink!55, fill=white, inner sep=3.05pt},
  ar/.style={-{Stealth[length=5.5pt]}, ink!65},
  leader/.style={quietink, line width=0.5pt},
]

% The box is the INK, not the ink plus every node's invisible inner sep: left is
% the axis line's left end, right is the right edge of the NIR label's glyphs,
% and the vertical pair reproduces the natural bounding box. 14.190cm in a
% 14.200cm block. Declared FIRST, so nothing drawn below can enlarge it; every
% mark in the picture sits inside it, so nothing spills into the margin either.
\path[use as bounding box] (0.713,-3.080) rectangle (14.903,1.290);

\node[lbl, anchor=south] at (3.28,0.74) {other drugs in the same comparison set};
\node[strong, anchor=south] at (10.07,0.74) {own truth};
% change (2): the leader. 0.35cm, quiet ink, label's south down to the dot's edge.
\draw[leader] (10.07,0.74) -- (10.07,0.39);

\fill[ink!10, rounded corners=1.5pt] (0.88,0.00) rectangle (10.07,0.44);
\foreach \x in {1.52, 2.80, 4.15, 5.67, 7.11, 8.71} {\node[other] at (\x,0.22) {};}
\node[own] at (10.07,0.22) {};
\node[other] at (11.43,0.22) {};
\node[other] at (12.54,0.22) {};
\node[anchor=west, font=\small] at (13.325,0.22) {$\NIR = \tfrac{6}{8}$};

\draw[ar] (0.72,-0.34) -- (13.02,-0.34);
\node[lbl, anchor=north] at (6.87,-0.48)
  {similarity of the prediction to each condition's truth $\longrightarrow$};

\draw[decorate, decoration={brace, mirror, amplitude=4.5pt}, ink!55]
  (0.88,-1.12) -- (10.07,-1.12);
\node[strong, anchor=north] at (5.43,-1.30) {the shaded fraction \emph{is} the score};

\node[lbl, anchor=west] at (0.72,-1.96) {own truth ranked first $\Rightarrow \NIR \to 1$};
\node[lbl, anchor=west] at (0.72,-2.38) {own truth ranked at random $\Rightarrow \NIR = 0.5$};
\node[strong, anchor=west] at (0.72,-2.80)
  {every similarity \emph{identical} $\Rightarrow \NIR = 0.5$ by the half-tie rule, not $0$};

\end{tikzpicture}%
\end{adjustbox}
% The comment character above is load-bearing. Without it the line break is an
% interword space inside the centred line, which spends 3.4pt of a measure the
% drawing now fills exactly and pushes the drawing 1.7pt off centre.
\caption[NIR is a discrimination score]{\captiontitle{$\NIR$ asks whether a prediction resembles its own
condition's truth more than it resembles the truths of the other conditions in the comparison set
$J$.} Own truth ranked first gives $\NIR \to 1$; ranked at random gives $0.5$; and a predictor whose
similarities are all identical --- the zero vector, which is what ``predict the average drug
response'' becomes in residual space --- scores $0.5$ by the half-tie rule, not $0$.}
\label{fig:nir}
\end{fullwidth}
\end{figure}
